\begin{frame}{템플릿 ②의 답안 뼈대}
  \begin{itemize}
    \item 시험에서 \textbf{``C를 크게 하면 무엇이 변하고 무엇이
      변하지 않는가''} --- 정답의 뼈대는 4구간표의 \textbf{구간
      IV}: \kq{``$C \ge C^*$ 이후 해는 $C$에 무관하다''}.
    \item 달성 가능한 마진의 범위는 \textbf{지오메트리가} 정한다
      --- 이상치를 무시하면 최대 3.0($-1$과 $+2$ 사이), 이상치를
      마진 안에 넣으려면 0.8($-1$과 $-0.2$ 사이)까지만 가능.
      C는 그 \textbf{두 극단 사이의 트레이드오프}를 선택할
      뿐, ``무한히 큰 마진''이라는 선택지는 존재하지 않는다.
    \item 실습 노트북 §2가 4구간표를 실제 $C$ 스위프로 검증
      ($\|w\|^2$가 $14.44C^2$로 오르면 4/9로 수평, $0.64C^2$
      로 다시 오르면 6.25로 수평), $\alpha_i$를
      \texttt{dual\_coef\_}로 읽어 손계산과 대조.
  \end{itemize}
\end{frame}
